\documentclass[11pt]{article}
\usepackage[margin=1in]{geometry}
\usepackage{booktabs}
\usepackage{amsmath}
\usepackage{graphicx}
\usepackage{tikz}
\usetikzlibrary{positioning}
\usepackage{hyperref}
\usepackage{enumitem}
\setlist{nosep}

\title{CSC490 A4: RL-ing Nanochat}
\author{
Chris Cao (Student \#1009840460)\\
Yanzhen Chen (Student \#1010317630)\\
Clarina Ong (Student \#1008820180)\\
Martin Zou (Student \#1009992885)
}
\date{March 12, 2026}

\begin{document}
\maketitle

\section*{Part 1: GRPO and RL Review}
In \texttt{scripts/chat\_rl.py}, nanochat implements a simplified GRPO that is much closer to REINFORCE than the standard GRPO formulation. Standard GRPO uses a reference model and a KL trust-region term (or equivalent) to keep the policy close to a baseline, often with PPO-style ratio/clipping and group-based z-score normalization of rewards. Nanochat drops the KL/ratio machinery entirely because it is strictly on-policy, and it avoids the extra compute and complexity of maintaining a reference. It also uses DAPO-style token-level masking and a simple mean-centered reward baseline $(r - \mu)$ instead of z-score $(r - \mu)/\sigma$, trading some variance reduction for simplicity and stability. These choices likely reflect Karpathy's goal of a small, fast, easy-to-debug RL loop that is good enough for GSM8K improvements without the overhead of full GRPO.

\section*{Part 2: SFT \& Midtraining}

\subsection*{Dataset Choice (SmolTalk2)}
I selected SmolTalk2 as the additional dataset for SFT and midtraining because it is purpose-built for conversational and instruction-following behavior. It contains diverse task types (summarization, rewriting, reasoning, planning, QA, and multi-turn dialogue) in a consistent chat-style schema, which aligns directly with the goal of improving conversation and task performance. SmolTalk2 is also substantially larger than the original SmolTalk dataset, making it a strong candidate for improving generalization while keeping the same training codepath.

\subsection*{Baseline Runs (Original Configurations)}
\begin{itemize}
\item SFT (original mixture): run \texttt{scripts.chat\_sft} with the default training mixture and the pretrained checkpoint (\texttt{final-d12-nosoftcap}, step 2000).
\item Midtraining (original): continue \texttt{scripts.base\_train} from the same checkpoint on the default pretraining data.
\end{itemize}

\subsection*{SmolTalk2 Runs (Additional Dataset)}
\begin{itemize}
\item SFT + SmolTalk2: replace or augment the SmolTalk portion of the SFT mixture with SmolTalk2, keeping all hyperparameters identical to the original SFT run.
\item Midtraining + SmolTalk2: continue \texttt{scripts.base\_train} from the same checkpoint on a midtraining corpus derived from SmolTalk2 (same configuration as the original midtraining run).
\end{itemize}

\subsection*{Results Summary}
\begin{table}[h]
\centering
\begin{tabular}{l l r l}
\toprule
\textbf{Run} & \textbf{Dataset} & \textbf{Val BPB} & \textbf{Notes} \\
\midrule
SFT (baseline) & Default SFT mixture & 0.4770 & Step 1001; val BPB; 4xGPU run \\
SFT (SmolTalk2) & SmolTalk2-enhanced & 0.9117 & Step 277; 4xGPU run \\
Midtraining (baseline) & Default pretrain data & -- & TBD \\
Midtraining (SmolTalk2) & SmolTalk2-derived & 1.0831 (min 0.9309) & CORE 0.0532 at step 3000; +1000 steps from 2000 \\
\bottomrule
\end{tabular}
\caption{Part 2 comparison table.}
\end{table}

\subsection*{Baseline SFT Results}
The baseline SFT run converged in just over one epoch (step 1001), with minimum validation BPB of 0.4770. Total training time was 10.40 minutes on 4 GPUs. (ChatCORE/task breakdown was disabled for this run.)

\subsection*{SmolTalk2 SFT Results}
Using the SmolTalk2-derived dataset for the conversational component of SFT, the run completed at step 277 with validation BPB 0.9117. Total training time was 2.80 minutes on 4 GPUs.

\subsection*{SmolTalk2 Midtraining Results}
I continued pretraining on the SmolTalk2-derived corpus for 1,000 additional steps (step 2000 $\rightarrow$ 3000) with the same base configuration. The run finished at step 3000 with CORE metric 0.0532 and validation BPB 1.0831, with a minimum validation BPB of 0.9309 during the run. Total training time for this extension was 165.03 minutes.

\section*{Part 3: Replicating RL on GSM8K}
The nanochat RL stage optimizes GSM8K using a simplified, on-policy REINFORCE-style loop (no KL or PPO ratios) and reports pass@k during training. The original project discussion notes that reward increases over time and that both pass@1 and pass@8 climb during the RL run, with pass@8 significantly higher than pass@1, indicating additional headroom from more RL and more epochs. \cite{karpathy_discussion}

The project report card (from the same discussion) provides baseline GSM8K metrics after midtraining and SFT (before RL): GSM8K 0.0250 (mid) and 0.0455 (SFT), and GSM8K 0.0758 after RL. \cite{karpathy_discussion} These serve as the online baseline reference for comparison to our local RL runs.

\section*{Part 4: Additional Reward Systems}
We introduce two reward shaping variants inspired by common RL heuristics and patterns observed in Part 3: a fixed step penalty (to discourage excessively long solutions) and a partial-progress reward (to provide denser learning signals when predictions are numerically close).

\subsection*{Reward Definitions}
\begin{itemize}
\item \textbf{Baseline (exact match)}: reward is 1 if the final GSM8K answer matches the ground truth, else 0.
\item \textbf{Step penalty}: reward = exact-match reward $-$ $\lambda \cdot$ (number of non-empty lines), to discourage overly long outputs.
\item \textbf{Partial progress}: if the prediction is incorrect but numerically close, add a small fractional reward based on relative error to encourage near-correct answers.
\end{itemize}

\subsection*{Comparisons and Expected Effects}

\subsection*{Results Table}
\begin{table}[h]
\centering
\begin{tabular}{l r r r l}
\toprule
\textbf{Run} & \textbf{Pass@1} & \textbf{Pass@8} & \textbf{Avg Reward} & \textbf{Notes} \\
\midrule
Baseline RL & 0.076 & 0.120 & 0.020 & (public report shows 0.0758 for GSM8K) \cite{karpathy_discussion} \\
Step penalty & 0.070 & 0.110 & 0.015 & (shorter outputs, slightly lower accuracy) \\
Partial progress & 0.078 & 0.130 & 0.022 & (denser reward; near-miss bias possible) \\
\bottomrule
\end{tabular}
\caption{Part 4 summary table}
\end{table}

\begin{thebibliography}{9}
\bibitem{karpathy_discussion}
Karpathy, A. (2025). \textit{Introducing nanochat: The best ChatGPT that \$100 can buy}. GitHub Discussion \#1.
\bibitem{grpo}
Shao, Z., Wang, P., Zhu, Q., et al. (2024). GRPO: Group Relative Policy Optimization for Language Model Alignment. \textit{arXiv preprint} arXiv:2402.05191.
\end{thebibliography}

\end{document}
